\documentclass[10pt,a4paper]{amsart}
\usepackage{amsmath,amssymb,amsthm,graphicx}
% Retained preamble for the Stacks Project chapter snapshot in this directory.
% The chapter's own preamble.tex is a copy of the shared Stacks Project preamble; the block
% below is sliced MECHANICALLY, line for line, from it. Nothing was added, removed or
% redefined; the wrapper class, the xr-hyper/hyperref wiring and the externaldocument list
% are omitted because the assembler supplies the document class.
%
% lines taken from preamble.tex: 3, 5, 9, 11, 12, 156, 157, 158, ... (51 lines total)
%
% package coverage restored from the chapter's own preamble (2026-10-13): xy, xy-2cell

\usepackage[all]{xy}
\xyoption{2cell}
\UseAllTwocells

\newenvironment{reference}{\comment}{\endcomment}
\newenvironment{slogan}{\comment}{\endcomment}
\newenvironment{history}{\comment}{\endcomment}
\theoremstyle{plain}
\newtheorem{theorem}[subsection]{Theorem}
\newtheorem{proposition}[subsection]{Proposition}
\newtheorem{lemma}[subsection]{Lemma}
\theoremstyle{definition}
\newtheorem{definition}[subsection]{Definition}
\newtheorem{example}[subsection]{Example}
\newtheorem{exercise}[subsection]{Exercise}
\newtheorem{situation}[subsection]{Situation}
\theoremstyle{remark}
\newtheorem{remark}[subsection]{Remark}
\newtheorem{remarks}[subsection]{Remarks}
\numberwithin{equation}{subsection}
\def\lim{\mathop{\mathrm{lim}}\nolimits}
\def\colim{\mathop{\mathrm{colim}}\nolimits}
\def\Spec{\mathop{\mathrm{Spec}}}
\def\Hom{\mathop{\mathrm{Hom}}\nolimits}
\def\Ext{\mathop{\mathrm{Ext}}\nolimits}
\def\SheafHom{\mathop{\mathcal{H}\!\mathit{om}}\nolimits}
\def\SheafExt{\mathop{\mathcal{E}\!\mathit{xt}}\nolimits}
\def\Sch{\mathit{Sch}}
\def\Mor{\mathop{\mathrm{Mor}}\nolimits}
\def\Ob{\mathop{\mathrm{Ob}}\nolimits}
\def\Sh{\mathop{\mathit{Sh}}\nolimits}
\def\NL{\mathop{N\!L}\nolimits}
\def\CH{\mathop{\mathrm{CH}}\nolimits}
\def\proetale{{pro\text{-}\acute{e}tale}}
\def\etale{{\acute{e}tale}}
\def\QCoh{\mathit{QCoh}}
\def\Ker{\mathop{\mathrm{Ker}}}
\def\Im{\mathop{\mathrm{Im}}}
\def\Coker{\mathop{\mathrm{Coker}}}
\def\Coim{\mathop{\mathrm{Coim}}}
\DeclareMathSymbol{\boxtimes}{\mathbin}{AMSa}{"02}
\def\QCohstack{\mathcal{QC}\!\mathit{oh}}
\def\Cohstack{\mathcal{C}\!\mathit{oh}}
\def\Spacesstack{\mathcal{S}\!\mathit{paces}}
\def\Quotfunctor{\mathrm{Quot}}
\def\Hilbfunctor{\mathrm{Hilb}}
\def\Curvesstack{\mathcal{C}\!\mathit{urves}}
\def\Polarizedstack{\mathcal{P}\!\mathit{olarized}}
\def\Complexesstack{\mathcal{C}\!\mathit{omplexes}}
\def\Pic{\mathop{\mathrm{Pic}}\nolimits}
\def\Picardstack{\mathcal{P}\!\mathit{ic}}
\def\Picardfunctor{\mathrm{Pic}}
\def\Deformationcategory{\mathcal{D}\!\mathit{ef}}
\usepackage{cleveref}
\crefname{originallemma}{original lemma}{original lemmas}
\Crefname{originallemma}{Original lemma}{Original lemmas}
\makeatletter
\newlabel{curves-lemma-rr}{{0BS6}{source}{}{sourcepointer-1}{}}
\newlabel{curves-lemma-rr@cref}{{[originallemma][0][]0BS6}{[1][0][]source}{}{}{}}
\makeatother
\crefname{originallemma}{original lemma}{original lemmas}
\Crefname{originallemma}{Original lemma}{Original lemmas}
\makeatletter
\newlabel{curves-lemma-tricanonical}{{0E8X}{source}{}{sourcepointer-2}{}}
\newlabel{curves-lemma-tricanonical@cref}{{[originallemma][0][]0E8X}{[1][0][]source}{}{}{}}
\makeatother
\crefname{originallemma}{original lemma}{original lemmas}
\Crefname{originallemma}{Original lemma}{Original lemmas}
\makeatletter
\newlabel{lemma-polarized-curves-in-polarized}{{0DPZ}{source}{}{sourcepointer-3}{}}
\newlabel{lemma-polarized-curves-in-polarized@cref}{{[originallemma][0][]0DPZ}{[1][0][]source}{}{}{}}
\makeatother
\crefname{originallemma}{original lemma}{original lemmas}
\Crefname{originallemma}{Original lemma}{Original lemmas}
\makeatletter
\newlabel{moduli-lemma-bounded-polarized}{{0E96}{source}{}{sourcepointer-4}{}}
\newlabel{moduli-lemma-bounded-polarized@cref}{{[originallemma][0][]0E96}{[1][0][]source}{}{}{}}
\makeatother
\crefname{originallemma}{original lemma}{original lemmas}
\Crefname{originallemma}{Original lemma}{Original lemmas}
\makeatletter
\newlabel{section-polarized-curves}{{0DPY}{source}{}{sourcepointer-5}{}}
\newlabel{section-polarized-curves@cref}{{[originallemma][0][]0DPY}{[1][0][]source}{}{}{}}
\makeatother
\title{Stacks, Moduli of Curves: lemma-stable-quasi-compact}
\author{The Stacks Project Authors}
\date{stacks-project:moduli-curves.tex:local-snapshot-20261012; GNU FDL 1.2 or later}
\begin{document}
\maketitle
\noindent\textit{Editorial note.} Statement span is lines 3503-3506 of sources/stacks-moduli-curves/pinned-original-source.tex; proof span is lines 3508-3554, the complete author proof (45 body lines).\par\medskip
\noindent\textit{Editorial note.} Navigation stubs: 5 cited-but-not-included labels are supplied as non-load-bearing pointers with their pinned tags where the snapshot tag index has them.\par\medskip
\noindent\textit{Editorial note.} extra\_packages is empty on purpose: the retained preamble already loads xy with options and a second declaration clashes in the assembler wrapper.\par\medskip
\noindent\textit{Editorial note.} Text is retained verbatim, including the author's own typos and placeholders.\par\medskip
\noindent\textit{Editorial note.} Queue index 12 of queue\_moduli-curves.json (pb=1491, stmt 3503-3506, proof 3508-3554).\par\medskip
\noindent\textit{Allowed standard prerequisites (editorial).}\begin{itemize}
\item Explicit prerequisite boundary. The cited inputs are prerequisites; the argument in the retained span is not replaced by any cited result.
\end{itemize}\medskip
\ifdefined\hypertarget\hypertarget{sourcepointer-1}{}\fi
\noindent\textit{Source reference.} Original lemma 0BS6: sources/stacks-moduli-curves/pinned-original-source.tex, source lines 3554--3554. Cited by the retained statement or proof but not part of the collected spans; supplied only as a navigation pointer. The argument does not depend on its content.\par\medskip
\ifdefined\hypertarget\hypertarget{sourcepointer-2}{}\fi
\noindent\textit{Source reference.} Original lemma 0E8X: sources/stacks-moduli-curves/pinned-original-source.tex, source lines 3554--3554. Cited by the retained statement or proof but not part of the collected spans; supplied only as a navigation pointer. The argument does not depend on its content.\par\medskip
\ifdefined\hypertarget\hypertarget{sourcepointer-3}{}\fi
\noindent\textit{Source reference.} Original lemma 0DPZ: sources/stacks-moduli-curves/pinned-original-source.tex, source lines 420--424. Cited by the retained statement or proof but not part of the collected spans; supplied only as a navigation pointer. The argument does not depend on its content.\par\medskip
\ifdefined\hypertarget\hypertarget{sourcepointer-4}{}\fi
\noindent\textit{Source reference.} Original lemma 0E96: sources/stacks-moduli-curves/pinned-original-source.tex, source lines 3554--3554. Cited by the retained statement or proof but not part of the collected spans; supplied only as a navigation pointer. The argument does not depend on its content.\par\medskip
\ifdefined\hypertarget\hypertarget{sourcepointer-5}{}\fi
\noindent\textit{Source reference.} Original lemma 0DPY: sources/stacks-moduli-curves/pinned-original-source.tex, source lines 370--402. Cited by the retained statement or proof but not part of the collected spans; supplied only as a navigation pointer. The argument does not depend on its content.\par\medskip
\begin{lemma}
\label{lemma-stable-quasi-compact}
Let $g \geq 2$. The stack $\overline{\mathcal{M}}_g$ is quasi-compact.
\end{lemma}

\begin{proof}
We will use the notation from Section \ref{section-polarized-curves}.
Consider the subset
$$
T \subset |\textit{PolarizedCurves}|
$$
of points $\xi$ such that there exists a field $k$ and a pair
$(X, \mathcal{L})$ over $k$ representing $\xi$
with the following two properties
\begin{enumerate}
\item $X$ is a stable genus $g$ curve, and
\item $\mathcal{L} = \omega_X^{\otimes 3}$.
\end{enumerate}
Clearly, under the continuous map
$$
|\textit{PolarizedCurves}|
\longrightarrow
|\Curvesstack|
$$
the image of the set $T$ is exactly the open subset
$$
|\overline{\mathcal{M}}_g| \subset |\Curvesstack|
$$
Thus it suffices to show that $T$ is quasi-compact.
By Lemma \ref{lemma-polarized-curves-in-polarized} we see that
$$
|\textit{PolarizedCurves}| \subset |\Polarizedstack|
$$
is an open and closed immersion. Thus it suffices to
prove quasi-compactness of $T$ as a subset of
$|\Polarizedstack|$. For this we use the criterion of
Moduli Stacks, Lemma \ref{moduli-lemma-bounded-polarized}.
First, we observe that for $(X, \mathcal{L})$
as above the Hilbert polynomial $P$ is the function
$P(t) = (6g - 6)t + (1 - g)$ by Riemann-Roch, see
Algebraic Curves, Lemma \ref{curves-lemma-rr}.
Next, we observe that $H^1(X, \mathcal{L}) = 0$
and $\mathcal{L}$ is very ample by
Algebraic Curves, Lemma \ref{curves-lemma-tricanonical}.
This means exactly that with $n = P(3) - 1$
there is a closed immersion
$$
i : X \longrightarrow \mathbf{P}^n_k
$$
such that $\mathcal{L} = i^*\mathcal{O}_{\mathbf{P}^1_k}(1)$
as desired.
\end{proof}

\end{document}
